%-----------------------------------------------------------------------
% Beginning of gsm-l-template.tex
%-----------------------------------------------------------------------
%
%    This is a template file for monographs to be published in the
%    Graduate Studies in Math series, for use with AMS-LaTeX 2.0.
%    Separate chapters should be included at the appropriate position.
%
%    Templates for various common text, math and figure elements are
%    given following the \end{document} line.
%
%    Remove any commented or uncommented macros you do not use.
%
%%%%%%%%%%%%%%%%%%%%%%%%%%%%%%%%%%%%%%%%%%%%%%%%%%%%%%%%%%%%%%%%%%%%%%%%

\documentclass{gsm-l}
%    For the GSM series, theorem heads are set in bold and not indented.
%    To use the standard "amsbook" style for theorem heads instead,
%    comment out the previous \documentclass line and use this one:
%\documentclass[theoremb]{gsm-l}

%    For use when working on individual chapters
%\includeonly{}

%    If you need symbols beyond the basic set, uncomment this command.
%\usepackage{amssymb}

%    If your book includes graphics, uncomment this command.
%\usepackage{graphicx}

%    If the book includes commutative diagrams, ...
%\usepackage[cmtip,all]{xy}

%    Include other referenced packages here.
\usepackage{}

\newtheorem{theorem}{Theorem}[chapter]
\newtheorem{lemma}[theorem]{Lemma}

\theoremstyle{definition}
\newtheorem{definition}[theorem]{Definition}
\newtheorem{example}[theorem]{Example}
\newtheorem{xca}[theorem]{Exercise}

\theoremstyle{remark}
\newtheorem{remark}[theorem]{Remark}

\numberwithin{section}{chapter}
\numberwithin{equation}{chapter}

%    For a single index; for multiple indexes, see the manual
%    "AMS Author Handbook, Monograph Classes", included in the
%    author package).
\makeindex

\begin{document}

\frontmatter

\title{}

%    Remove any unused author tags.

%    author one information
\author{}
\address{}
\curraddr{}
\email{}
\thanks{}

%    author two information
\author{}
\address{}
\curraddr{}
\email{}
\thanks{}

%    The 2020 edition of the Mathematics Subject Classification is
%    the current definitive version.
\subjclass[2020]{Primary }

\keywords{}

\maketitle

%    Dedication.  If the dedication is longer than a few lines,
%    remove the centering instructions and the line break.
%\cleardoublepage
%\thispagestyle{empty}
%\vspace*{13.5pc}
%\begin{center}
%  Dedication text (use \\[2pt] for line break if necessary)
%\end{center}
%\cleardoublepage

%    Change page number to 7 if a dedication is present.
\setcounter{page}{5}

\tableofcontents

%    Include unnumbered chapters (preface, acknowledgments, etc.) here.
\include{}

\mainmatter
%    Include main chapters here.
\include{}

\appendix
%    Include appendix "chapters" here.
\include{}

\backmatter
%    Bibliographies can be prepared with BibTeX using amsplain,
%    amsalpha, or (for author-year style) natbib.
\bibliographystyle{amsalpha}
\bibliography{}

%    See note above about multiple indexes.
\printindex

\end{document}

%%%%%%%%%%%%%%%%%%%%%%%%%%%%%%%%%%%%%%%%%%%%%%%%%%%%%%%%%%%%%%%%%%%%%%%%

%    Templates for common elements of a mathematics book; for additional
%    information, see the AMS Author Handbook, Monograph Classes,
%    included in every AMS author package, and the amsthm user's guide,
%    linked from http://www.ams.org/tex/amslatex.html .

%    Section headings
\section{}
\subsection{}

%    Ordinary theorem and proof
\begin{theorem}[Optional addition to theorem head]
% text of theorem
\end{theorem}

\begin{proof}[Optional replacement proof heading]
% text of proof
\end{proof}

%    Figure insertion; default placement is top; if the figure occupies
%    more than 75% of a page, the [p] option should be specified.
\begin{figure}
\includegraphics{filename}
\caption{text of caption}
\label{}
\end{figure}

%    Mathematical displays; for additional information, see the amsmath
%    user's guide:  texdoc amamath  or
%  http://mirror.ctan.org/tex-archive/macros/latex/required/amsmath/amsldoc.pdf

% Numbered equation
\begin{equation}
\end{equation}

% Unnumbered equation
\begin{equation*}
\end{equation*}

% Aligned equations
\begin{align}
  &  \\
  &
\end{align}

%-----------------------------------------------------------------------
% End of gsm-l-template.tex
%-----------------------------------------------------------------------
